\chapter{The basepoint one}\label{chap:basepoint}

Section~\ref{ch:deconv} recovered the zeros from the Thorin density at the resolution $\varepsilon=0.1$, close to the critical line, where the positivity of the density at every height is open. We now work at the resolution $\varepsilon=\frac12$ of the Cauchy flow $\varepsilon\mapsto\rho_\varepsilon$ of Section~\ref{ch:clock}, the basepoint one. The completed zeta function has two probabilistic readings. On the Mellin side, Williams \cite{Williams90} and Biane, Pitman and Yor \cite{BPY01} showed that $2\xi(s)$ is the Mellin transform of a sum of two Bessel hitting times, which gives the P\'olya variable and its exponential tilts at every basepoint with no hypothesis. On the Laplace side, the shifted clocks $T_\alpha$, with $\E e^{-uT_\alpha}=\xi(\alpha)/\xi(\alpha+\sqrt u)$, carry the Thorin programme. At the centre, $\alpha=\frac12$, the shifted clock is the clock of \cite{P17}, which is infinitely divisible and self decomposable with no hypothesis (Theorem~\ref{thm:ch12:idsd}) and a generalised gamma convolution exactly when RH holds (Proposition~\ref{prop:rigid}). At a basepoint $\alpha$ to the right of every zero, $T_\alpha$ is a GGC whose Thorin density is the Poisson smoothing of the zeros by Cauchy kernels of widths $\alpha-\beta$ (Theorem~\ref{thm:poisson} and Proposition~\ref{prop:struct}), and from $\alpha=1$ on this density is positive with no hypothesis, because no zero lies on the line $\re s=1$ or to its right. The basepoint one is therefore the leftmost basepoint at which positivity is a theorem, and the natural place to separate what is unconditional from what is equivalent to RH.

At the basepoint one everything can be computed and most things can be proved, and we work the example out completely, as a test bed for the tools of the book. We ask three questions of $T_1$: what its Thorin measure and its law are, whether its Thorin density comes from the sine squared lemma, and whether it is a hitting time to which Kent's theorem applies. Section~\ref{ch:basepoint} answers the first two: it computes the Thorin measure and the law of $T_1$, and obtains the Thorin density from a L\'evy form of $\log\xi(1+s)$ by the sine squared lemma. It also proves, with no hypothesis and with a computation of the zeros below height $2000$ in ball arithmetic and explicit zero counts above it, that the density is smallest at the origin, so that the drift of the L\'evy form becomes a positive stable one half summand of $T_1$, independent of a GGC remainder (Proposition~\ref{bp:prop:v1min}). Below the basepoint one it proves the same on the verified window for every $\alpha\in(\frac12,7.35]$, and shows that the minimum at the origin at every height for every $\alpha\in(\frac12,1)$ is equivalent to RH (Section~\ref{sec:dbeta:comp}). Section~\ref{ch:bphit} answers the third: every $T_\alpha$ with $\alpha>\frac12$ is a hitting time, the tilted centre clock run through the inverse local time of a sticky Brownian motion, with an exact Brownian boundary that, as computed, becomes flat in the right tail.

For the Thorin measure and the L\'evy form alike, what is unconditional at the basepoint one extends to every basepoint to the right of the line $\re s=\alpha$, for a given $\alpha<1$, exactly when no zero lies to the right of that line; required for every $\alpha>\frac12$, this is RH. Chapters~\ref{chap:strip} and \ref{chap:almost} carry these objects across the critical strip and read them in other settings, and Section~\ref{ch:bpeta} repeats the calculations of this chapter for the function $L(s,\chi_{12})$ of Section~\ref{ch:chi12}, whose theta series is the Dedekind eta function.
